% ============================================================
%  CAPÍTULO 4: DATOS Y DESCRIPCIÓN EMPÍRICA
% ============================================================
\chapter{Datos y descripción empírica}
\label{cap:datos}

% ------------------------------------------------------------
\section{Introducción}
% ------------------------------------------------------------

El análisis empírico de este trabajo descansa sobre un panel de datos de
elaboración propia, construido a partir de la compilación y homogeneización de
tres fuentes estadísticas oficiales. El panel cubre las 17 comunidades autónomas
españolas durante el período 2016--2024, con 153 filas teóricas (17 CCAA $	imes$
9 años). Tras eliminar valores ausentes: (i) las variables de gasto están
disponibles en 150 observaciones, (ii) el PIB per cápita en 136 y (iii) las
listas de espera en 153 (todas las CCAA). Al construir retardos, la muestra
efectiva se reduce a 118 observaciones (17 CCAA) para el modelo base con un
retardo y a 101 (17 CCAA) cuando se utilizan dos retardos. Esta sección describe las fuentes utilizadas, el
proceso de construcción del panel y las principales características de las
variables empleadas.

% ------------------------------------------------------------
\section{Fuentes de datos}
% ------------------------------------------------------------

\subsection{Gasto en seguros sanitarios privados --- Encuesta de Presupuestos
Familiares (INE)}

La variable dependiente se obtiene de la \textbf{Encuesta de Presupuestos
Familiares (EPF)} del Instituto Nacional de Estadística \parencite{INE_EPF}. La
EPF recoge el gasto monetario de los hogares residentes en España, desagregado por
comunidades autónomas y clasificado según la nomenclatura COICOP. La partida
utilizada corresponde al epígrafe \textbf{06.3: Servicios hospitalarios y de
seguros sanitarios privados}.

Se utiliza el gasto por persona protegida (€/p.p.), expresado en términos
\textbf{reales} mediante deflactación por el Índice de Precios al Consumo
correspondiente a la rúbrica de seguros sanitarios (IPC seguros, base 2025 = 100):

\begin{equation}
  \text{GastoReal}_{it} =
  \frac{\text{GastoNominal}_{it}}{\text{IPC\_Seguros}_{it}} \times 100
  \label{eq:deflactacion}
\end{equation}

\subsection{Producto Interior Bruto per cápita --- Contabilidad Regional de España
(INE)}

El PIB per cápita a precios constantes (base 2015) se obtiene de la
\textbf{Contabilidad Regional de España (CRE)} \parencite{INE_CRE}. El PIB real
per cápita se construye dividiendo el PIB a precios constantes de cada CCAA entre
su población de referencia a 1 de enero de cada año (Padrón Municipal de
Habitantes). Esta variable actúa como \textit{proxy} de la \textbf{renta
disponible agregada} de cada territorio.

\subsection{Tiempos de espera en consulta de especialistas --- SISLE-SNS
(Ministerio de Sanidad)}

Los datos sobre tiempos de espera provienen del \textbf{Sistema de Información
sobre Listas de Espera del SNS (SISLE-SNS)}, gestionado por el Ministerio de
Sanidad \parencite{MinSanidad}. Este registro publica semestralmente los días
medios de espera para consulta con especialista, desagregados por comunidad
autónoma y especialidad médica.

La variable utilizada es la \textbf{espera media en consulta de especialistas}
(días), justificada teóricamente en el Capítulo~\ref{cap:teoria} mediante la Ley
de Little. Para las estimaciones se utilizan los datos de cierre de año
(diciembre) y se construyen los retardos temporales (lag 1 y lag 2) dentro de
cada CCAA.

\subsection{Variables demográficas y de control}

El porcentaje de población mayor de 65 años se obtiene del \textbf{Padrón
Continuo de Población} \parencite{INE_CRE}, que proporciona series anuales de
estructura por edad para cada comunidad autónoma. La variable \textit{dummy} de
pandemia COVID-19 toma valor 1 para el año 2020 en todas las CCAA.

% ------------------------------------------------------------
\section{Estadísticos descriptivos}
% ------------------------------------------------------------

La Tabla~\ref{tab:descriptivos} presenta los estadísticos descriptivos de las
variables principales del análisis. El panel bruto contiene 153 observaciones
(17 CCAA $\times$ 9 años), pero el número de observaciones válidas varía por
variable (ver nota al pie) y los modelos econométricos usan subconjuntos entre
118 y 101 observaciones tras construir retardos.

\begin{table}[htbp]
  \centering
  \caption{Estadísticos descriptivos del panel (17\,CCAA,\,2016--2024)}
  \label{tab:descriptivos}
  \begin{threeparttable}
    \small
    \resizebox{\textwidth}{!}{%
    \begin{tabular}{llS[table-format=6.2]S[table-format=5.2]
                       S[table-format=5.2]S[table-format=5.2]S[table-format=6.2]}
      \toprule
      \textbf{Variable}
        & \textbf{Descripción}
        & {\textbf{Media}}
        & {\textbf{D.\,Típica}}
        & {\textbf{Mínimo}}
        & {\textbf{Mediana}}
        & {\textbf{Máximo}} \\
      \midrule
      \texttt{g\_seguros\_real}
        & Gasto real seguros (€/p.p.)
        & 237.70 & 162.06 & 57.26 & 176.83 & 779.37 \\
      \texttt{log\_gasto\_real}
        & Log del gasto real
        & 5.277 & 0.606 & 4.048 & 5.175 & 6.658 \\
      \texttt{pib\_pc}
        & PIB per cápita real (€, p.c. 2015)
        & 25996 & 5426 & 17239 & 24782 & 42639 \\
      \texttt{log\_pib\_pc}
        & Log del PIB per cápita
        & 10.145 & 0.204 & 9.755 & 10.118 & 10.661 \\
      \texttt{esp\_esp}
        & Espera media especialista (días)
        & 78.3 & 32.6 & 27 & 73 & 176 \\
      \texttt{pob\_65\_pct}
        & Población $\geq 65$ años (\%)
        & 20.48 & 3.21 & 14.2 & 20.8 & 26.7 \\
      \bottomrule
    \end{tabular}}% fin resizebox
    \begin{tablenotes}[flushleft]
      \small
      \item Fuentes: INE (EPF, CRE, Padrón); Ministerio de Sanidad (SISLE-SNS).
        Gasto deflactado por IPC Seguros
        (base 2025 = 100).
      \item Observaciones válidas: 150 para gasto real en seguros y su logaritmo;
        136 para PIB per cápita y su logaritmo; 153 para espera media en
        especialistas; y 153 para porcentaje de población de 65 o más años.
      \item Las variables binarias de pandemia (año 2020) y espera alta
        (espera superior a 100 días) no se incluyen en esta tabla descriptiva.
    \end{tablenotes}
  \end{threeparttable}
\end{table}

\subsection{Heterogeneidad transversal del gasto}

El gasto real en seguros privados presenta una \textbf{dispersión muy elevada}
entre comunidades autónomas: la desviación típica (162~€/p.p.) representa el 68\%
de la media (238~€/p.p.). El coeficiente de variación implica que las CCAA con
mayor gasto privado (Madrid, País Vasco, Cataluña) gastan entre 4 y 6 veces más
en seguros privados que las de menor propensión (Extremadura, Murcia, Canarias en
algunos años). El rango observado ---de 57 a 779~€/p.p.--- subraya que la sanidad
privada en España no es un fenómeno homogéneo sino profundamente territorial.

\subsection{Capacidad explicativa del PIB per cápita}

El PIB per cápita oscila entre 17.239~€ y 42.639~€, con una media de 25.996~€.
La correlación simple entre $\log(\text{gasto\_real})$ y $\log(\text{pib\_pc})$
alcanza $r = 0{,}625$ ($p < 0{,}001$, $N = 134$). Una regresión MCO simple entre
ambas variables (134 observaciones válidas) produce una pendiente de 1{,}87 con
$R^2 = 0{,}39$, sugiriendo que la renta per cápita explica por sí sola casi el
40\% de la variación total en el gasto privado.

La diferencia respecto a las 136 observaciones disponibles para
	exttt{log\_pib\_pc} se debe a que el cruce bivariante exige datos completos
simultáneos de $\log(\text{pib\_pc})$ y $\log(\text{gasto\_real})$: hay 2 casos
con PIB observado pero gasto faltante, por lo que la muestra conjunta se reduce a
$N = 134$.

\subsection{Variabilidad de las listas de espera}

Los tiempos de espera exhiben la mayor variabilidad relativa del panel: la
desviación típica (32{,}6 días) representa el 42\% de la media (78{,}3 días). El
rango va desde 27 hasta 176 días. El umbral de 100 días ---adoptado como criterio
de espera excesiva en el modelo M3--- es superado en 33 observaciones (22\% del
total). La baja correlación simple entre $\texttt{esp\_esp}$ y
$\log(\text{gasto\_real})$ ($r = 0{,}031$) no debe interpretarse como ausencia de
efecto, dado que el mecanismo es \textbf{diferido}.

\subsection{Envejecimiento demográfico}

El porcentaje de población mayor de 65 años oscila entre el 14{,}2\% y el 26{,}7\%,
con una media del 20{,}5\%. La correlación negativa con el gasto en seguros
privados ($r = -0{,}35$) sugiere que el envejecimiento no estimula el gasto
privado, posiblemente porque los mayores dependen más intensamente del sistema
público y tienen menor capacidad de pago para primas de seguro.

% ------------------------------------------------------------
\section{Evolución temporal y correlación entre variables clave}
% ------------------------------------------------------------

La Figura~\ref{fig:scatter} presenta el diagrama de dispersión entre el logaritmo
del PIB per cápita y el logaritmo del gasto real en seguros privados para las
observaciones del panel con datos completos en ambas variables ($N = 134$).
Para este cruce concreto, el período efectivo es 2016--2023, ya que en 2024 no
hay observaciones con PIB per cápita disponible.

\begin{figure}[htbp]
  \centering
  \includegraphics[width=0.85\textwidth]{scatter_pib_gasto}
  \caption[Gasto real en seguros privados y PIB per cápita]{%
    Gasto real en seguros privados y PIB per cápita
    (panel efectivo 2016--2023).}
  \label{fig:scatter}
  \begin{minipage}{0.85\textwidth}
    \footnotesize\vspace{4pt}
    \textit{Nota:} Cada punto representa una observación CCAA-año con datos
    completos en $\log(\text{pib\_pc})$ y $\log(\text{gasto\_real})$
    ($N = 134$). Línea de tendencia MCO: $\hat{\beta} = 1{,}87$,
    $R^2 = 0{,}39$, $p < 0{,}001$. Fuentes: INE (EPF, CRE).
  \end{minipage}
\end{figure}

La figura revela tres patrones relevantes. \textbf{Primero}, existe una relación
positiva y estadísticamente significativa entre renta y gasto en seguros privados.
\textbf{Segundo}, la nube de puntos muestra heterogeneidad transversal persistente
que los efectos fijos capturan. \textbf{Tercero}, las observaciones más recientes
tienden a situarse ligeramente por encima de los años anteriores, coherente con la
tendencia creciente del gasto en seguros privados.

% ------------------------------------------------------------
\section{Restricciones de los datos y consideraciones para el análisis}
% ------------------------------------------------------------

El panel presenta las siguientes restricciones:

\begin{itemize}
  \item \textbf{Valores ausentes:} Tras eliminar valores perdidos, la muestra
    efectiva es de 150 observaciones para las variables de gasto, 136 para el PIB
    per cápita y 153 para listas de espera (las 17 CCAA disponen de datos de
    espera para todo el período). Al construir retardos, las estimaciones
    principales usan \textbf{118 observaciones} (modelo con un retardo, 17 CCAA)
    y \textbf{101 observaciones} (modelo con dos retardos, 17 CCAA).
  \item \textbf{Retardos y período efectivo:} La construcción de retardos de tiempo
    de espera (lag 1 y lag 2) reduce el período de estimación a 2017--2024 para los
    modelos con un retardo y a 2018--2024 para los modelos con dos retardos.
  \item \textbf{Ámbito territorial:} El panel incluye las 17 comunidades
    autónomas españolas (peninsulares e insulares). Ceuta y Melilla no forman
    parte del ámbito del estudio al no ser comunidades autónomas.
\end{itemize}
